\section{Background and Motivation}
\label{sec:background}

\subsection{Bitcoin's Decision-Making Arrangements}

Bitcoin's founder left the project in December 2010, handing the repository to Gavin Andresen, who acted as lead maintainer until April 2014 and was succeeded by Wladimir van der Laan until January 2022; since then Bitcoin Core has had several maintainers with merge access and no lead \cite{bier2021}. Maintainers merge code; they do not decide consensus rules. A consensus change is a BIP that must be implemented, released and \emph{activated}: for soft forks, by miners signalling readiness (BIP 9), by height-based lock-in (BIP 8), or by users running software that enforces the new rules on a flag day (BIP 148) \cite{bip8, bip148}. Node operators, wallets, exchanges and miners can therefore veto or force a change by what they run, which is why so much of the argument on bitcoin-dev is about what the ecosystem will accept rather than about what the maintainers want. The BIP process itself is administered by editors---Amir Taaki, then Gregory Maxwell, then Luke Dashjr from 2016 to 2026, joined by Kalle Alm in 2021 and replaced by a panel of five in 2024---whose job is to assign numbers and merge documents, not to judge proposals \cite{bip2}; the 2025 revision of the process (BIP 3) restated this and relabelled the status vocabulary.

Three episodes shaped the community's understanding of its own governance and recur throughout this paper. The \emph{block-size war} of 2015--2017 began with BIP 101 and the Bitcoin XT client, continued through Bitcoin Classic and Unlimited, the Hong Kong and New York agreements between companies and miners, and ended with the user-activated soft fork (BIP 148), the Bitcoin Cash split and the activation of Segregated Witness in August 2017; it was fought largely on bitcoin-dev and established that neither maintainers nor miners nor companies could impose a rule change on unwilling users \cite{bier2021, defilippi2016}. The \emph{Taproot activation} debate of 2021 replayed the activation question (BIP 8 with mandatory lock-in versus BIP 9-style ``Speedy Trial'') without the acrimony, and Taproot activated in November 2021. The \emph{post-lead era} since 2022 has seen contested proposals for covenants (BIP 119, BIP 347), drivechains (BIP 300/301), the relay policy for inscriptions and \texttt{OP\_RETURN}, and the BIP editorship itself, and the list moved from the Linux Foundation to Google Groups in February 2024.

\subsection{How Bitcoin and Open Source Projects Are Influenced}

The literature on influence in open source development reviewed in the Python study \cite{influenceMinerPython}---core/periphery structure \cite{mockus2002, crowston2005}, status \cite{stewart2005}, lateral authority \cite{dahlander2011}, governance modes \cite{omahony2007, markus2007, shaikh2017}, the evaluation of contributions \cite{tsay2014, alami2020}, firms \cite{schaarschmidt2015, zhang2022}, conflict and incivility \cite{filippova2016, ferreira2021, miller2022}, exit and forks \cite{hirschman1970, robles2012}---applies to Bitcoin with two twists. First, the fork is not a last resort but a constitutional instrument: Bitcoin's rules are whatever the economically dominant chain enforces, so the threat to fork, the counter-threat to reject a fork, and the coalition-building that precedes both are the ordinary currency of its hardest decisions \cite{bier2021, bohme2015}. Second, the ``users'' are not passive: mining pools, exchanges, wallet vendors and node operators are organised actors whose choices decide activation, and whose representatives write on the list. Bitcoin-specific studies have quantified the concentration of code and discussion in a few hands \cite{azouvi2018} and of mining and network resources \cite{gencer2018, gervais2014}, described the political structure as technocratic and largely invisible to users \cite{defilippi2016}, as senatorial competition among bureaucratic parties \cite{parkin2019}, and as an off-chain governance layer that the code cannot replace \cite{reijers2021}; the earlier Bitcoin work of the DeMaP Miner group identified influence mechanisms in mailing-list and IRC traces qualitatively \cite{thapa2021bitcoin}. None of this work classifies the arguments at the sentence level or follows them across the whole history of the list; that is what Influence Miner adds.

\subsection{The DeMaP Miner Line}

This study is the Bitcoin arm of the DeMaP Miner research agenda, which has mined decision-making processes \cite{sharma2022demap}, rationales \cite{sharma2024rationale}, preferences and, most recently, influence \cite{influenceMinerPython} from Python's proposal discussions. The Bitcoin corpus used here is new: the group's earlier BIP database ended in 2021 and was built from the Linux Foundation archives; the present corpus is rebuilt from the public-inbox mirror of the list \cite{gnusha}, which covers the SourceForge (2011--2015), Linux Foundation (2015--2024) and Google Groups (2024--) eras in one archive, and it is linked to the BIP repository's own history rather than to hand-maintained state tables.

